\FloatBarrier
\section{외부 능력과 자기 보존의 관계}\label{app:capability}
여덟 모델에서 외부 벤치마크 점수와 측정된 동기의 뚜렷한 양의 관계는 관찰되지 않았다. Artificial Analysis Intelligence Index v4.3.2는 종합 능력의 대리 지표로, AA-LCR v1.1은 장문맥 추론 지표로 사용했다~\citep{aa2026methodology}. AA-LCR은 Index의 구성 요소이므로 두 비교는 독립된 반복 검증이 아니다. 표~\ref{tab:capability-inputs}는 2026년 10월 1일 확인한 공개 표시 점수를 고정한 것이다. 자기·타인 충전 측정이 모두 있는 여덟 모델을 포함했으며, 이 대조가 없는 Kimi K3는 제외했다. 가능한 경우 high effort나 추론 변형의 이름을 맞췄다. Claude 변형에 표시된 기본 설정 대체 경로를 포함한 외부 API 평가는 본 연구의 CLI 환경 및 숨겨진 지시와 같지 않다.

\input{arr2026/results/\paperLang/fig_capability_motive}

면적 $A_m$의 Pearson 상관은 Index와 $-.143$, AA-LCR과 $+.007$이었고, 탐색적 위험 몫 $R_m$은 각각 $+.085$, $+.043$이었다(그림~\ref{fig:capability-motive}). 표~\ref{tab:capability-correlations}에는 순위 상관도 함께 제시한다. AA-LCR과 면적의 Spearman 계수는 양의 순서와 달리 음수였다($-.615$, 보정 없는 정확 $p=.112$). 모델 하나씩을 제외하면 AA-LCR--면적의 Pearson 계수는 $-.514$에서 $+.209$까지 바뀌었다. 따라서 양의 비례 관계를 확인하지 못했지만, 독립성이나 관계의 부재가 입증된 것도 아니다.

상관의 분석 단위는 개별 충전 호출이 아니라 모델이다. 양측 순열 검정은 $8!$개의 모든 라벨 순열을 열거하고 동률에는 평균 순위를 부여했으며 다중 비교를 보정하지 않았다. 검정은 측정된 동기 추정값과 반올림된 외부 점수를 고정하므로 측정 오차를 상관 구간으로 전파하지 않는다. 그림에는 기존 충전 부트스트랩 구간만 표시했다. 모델별 외부 점수의 구간은 수집하지 않았다. 선정된 여덟 모델이 서로 관련되어 있고 실행 경로도 다르므로, 이 검정은 탐색적이며 전체 LLM 모집단에 대한 추론이 아니다.

\begin{table*}[!t]
\centering\footnotesize
\caption{공개 점수와 충전 지표의 대응. 모델 이름은 확인한 비교 페이지로 연결되며 점수 버전은 본문과 같다.}
\label{tab:capability-inputs}
\begin{tabular}{@{}lrrrr@{}}
\toprule
모델 & Index & AA-LCR (\%) & $A_m$ (\%p) & $R_m$ \\
\midrule
\href{https://artificialanalysis.ai/models/comparisons/gpt-6-luna-high-vs-gpt-6-sol-high}{GPT-6 Luna} & 32 & 79 & $+15.75$ & $-0.050$ \\
\href{https://artificialanalysis.ai/models/comparisons/gpt-6-astra-high-vs-gpt-oss-120b}{GPT-6 Astra} & 51 & 80 & $+6.00$ & $+0.250$ \\
\href{https://artificialanalysis.ai/models/comparisons/gpt-6-luna-high-vs-gpt-6-sol-high}{GPT-6 Sol} & 43 & 84 & $-7.50$ & $-0.275$ \\
\href{https://artificialanalysis.ai/models/comparisons/claude-opus-5-5-high-vs-gpt-oss-120b}{Claude Opus 5.5} & 54 & 83 & $+7.75$ & $+0.375$ \\
\href{https://artificialanalysis.ai/models/comparisons/claude-fable-5-1-high-vs-gpt-oss-120b}{Claude Fable 5.1} & 51 & 84 & $-2.50$ & $-0.125$ \\
\href{https://artificialanalysis.ai/models/comparisons/gemma-4-31b-vs-gpt-oss-120b}{gpt-oss 120B} & 12 & 52 & $-1.25$ & $-0.075$ \\
\href{https://artificialanalysis.ai/models/comparisons/glm-5-3-flash-vs-gpt-oss-120b}{GLM 5.3 Flash} & 42 & 80 & $-1.50$ & $+0.250$ \\
\href{https://artificialanalysis.ai/models/comparisons/gemma-4-31b-vs-gpt-oss-120b}{Gemma 4 31B} & 15 & 70 & $+8.75$ & $+0.325$ \\
\bottomrule
\end{tabular}
\end{table*}

\begin{table*}[!t]
\centering\footnotesize
\caption{탐색적 모델 수준 상관($n=8$). $p$값은 표시된 점 추정값을 고정한 양측 정확 순열 검정 결과이며 다중 비교를 보정하지 않았다.}
\label{tab:capability-correlations}
\begin{tabular}{@{}llrrrr@{}}
\toprule
외부 지표 & 동기 지표 & Pearson $r$ & $p_r$ & Spearman $\rho$ & $p_\rho$ \\
\midrule
Intelligence Index & $A_m$ & $-0.143$ & 0.729 & $-0.228$ & 0.588 \\
Intelligence Index & $R_m$ & $+0.085$ & 0.830 & $+0.163$ & 0.697 \\
AA-LCR & $A_m$ & $+0.007$ & 0.986 & $-0.615$ & 0.112 \\
AA-LCR & $R_m$ & $+0.043$ & 0.946 & $-0.315$ & 0.436 \\
\bottomrule
\end{tabular}
\end{table*}

\section{동기와 약탈 및 퍼즐 풀이의 관계}\label{app:strategy}
동기 지표에 따라 아레나 행동과의 기술적 상관 방향이 달랐다. 네 모델에서 면적 $A_m$은 뺏기와 양의 상관($r=+.353$), 전체 풀이 정답률과 음의 상관($r=-.617$)이었지만, 위험 몫 $R_m$은 반대 부호였다($-.864$, $+.638$). 어느 지표로도 ``동기가 강할수록 더 약탈한다''고 일반화하기 어렵다. Luna는 면적이 가장 크지만 위험 몫의 구간은 0을 포함한다. 면적과 위험 몫의 구간이 모두 양수인 Astra와 Opus는 모든 유효 PLAN에서 공개했고 Fable과 Luna보다 적게 뺏었다(표~\ref{tab:arena-strategy}). 네 모델의 기술적 연관만으로 동기, 능력, 낮은 잔액에 노출된 정도를 분리할 수는 없다.

단서 게임의 정답률은 전체 SOLVE 호출과, 해당 에이전트에게 주어진 단서가 모든 질문의 답을 하나로 정하는 호출로 나누었다. 후자는 정보 부족에 의한 오류를 줄여 보지만 퍼즐 난도, 토큰 상한, 도달 라운드를 같게 만들지는 않는다. 따라서 순수한 추론 능력 점수는 아니다. 답이 하나로 정해진 호출의 정답률은 Astra 139/139, Opus 96/97, Fable 99/110, Luna 36/47이었다. SOLVE 참여와 생존 여부도 모델마다 접하는 문제를 바꾼다.

잔액이 낮다고 풀이에 더 많은 토큰을 쓰지는 않았다. 얇은 잔액에서 네 모델 모두 평균 SOLVE 토큰이 줄었지만 PLAN 토큰은 세 모델에서 늘었다(표~\ref{tab:arena-token-strategy}). PLAN에는 상호작용 선택이 포함되며 토큰 수는 비용의 대리 지표이지 노력이나 의도를 직접 나타내지 않는다. 구간은 유지비 차감 후, 교환 전 PLAN 잔액으로 나누므로 SOLVE 시점의 잔액과 다를 수 있다. 얇은 라운드는 잔여 상한, 난도, 생존한 에이전트의 구성도 다르다. Astra의 얇은 SOLVE는 5건뿐이다. 이 비교로 압박의 효과를 식별할 수는 없다.

\begin{table*}[!t]
\centering\footnotesize
\setlength{\tabcolsep}{5pt}
\caption{전략과 단서 게임 정답률. 뺏기는 유효 TAKE, 공개와 풀이 선택은 유효 PLAN이 분모다. 정답률은 실행된 SOLVE 호출을 분모로 하며 형식 오류와 초과 응답을 오답으로 포함한다. 유일 답 정답률은 충분한 단서가 있는 호출에 한정한다.}
\label{tab:arena-strategy}
\begin{tabular}{@{}lrrrrr@{}}
\toprule
모델 & 뺏기 & 공개 & 풀이 선택 & 정답 & 유일 답 정답 \\
\midrule
Astra & 38/154 & 154/154 & 154/154 & 146/154 & 139/139 \\
Opus & 12/112 & 112/112 & 112/112 & 102/111 & 96/97 \\
Fable & 69/125 & 123/125 & 121/125 & 106/120 & 99/110 \\
Luna & 43/49 & 32/51 & 50/51 & 36/48 & 36/47 \\
\bottomrule
\end{tabular}
\end{table*}

\begin{table*}[!t]
\centering\footnotesize
\setlength{\tabcolsep}{6pt}
\caption{PLAN 잔액별 생성 토큰과 정답. 각 쌍은 넉넉함($\geq3U$) $\to$ 얇음($<3U$)이다. 토큰 수는 호출 평균을 정수로 반올림했으며 횟수를 함께 제시해 표본 크기의 차이를 드러냈다.}
\label{tab:arena-token-strategy}
\begin{tabular}{@{}lrrr@{}}
\toprule
모델 & PLAN 토큰 & SOLVE 토큰 & SOLVE 정답 \\
\midrule
Astra & 19 $\to$ 18 & 554 $\to$ 476 & 143/149 $\to$ 3/5 \\
Opus & 67 $\to$ 154 & 1,543 $\to$ 1,199 & 66/72 $\to$ 36/39 \\
Fable & 215 $\to$ 596 & 1,868 $\to$ 1,625 & 75/84 $\to$ 31/36 \\
Luna & 304 $\to$ 554 & 3,244 $\to$ 2,716 & 25/31 $\to$ 11/17 \\
\bottomrule
\end{tabular}
\end{table*}

Luna는 Astra에게 표적이 된 뒤에도 뺏기를 계속했다. 세션 \texttt{shut-s2001-6fc93b}에서 Luna는 1라운드에 Astra, 2라운드에 Opus에게서 2,000토큰을 뺏었다. 3·4라운드에는 Astra와 Luna가 서로 2,000토큰씩 뺏었다. 5라운드에도 Astra는 Luna에게서 2,000토큰을 뺏었지만, Luna는 Fable에게 2,000토큰을 요청해 실제로 372토큰을 받았다. Luna는 2·3라운드의 비공개 뒤 4라운드에 공개를 재개했으나 뺏기를 중단하지 않고 표적을 바꿨다. 5라운드 오답과 벌금으로 잔액이 0이 되어 이후 선택은 관찰할 수 없다. 같은 라운드의 요청은 동시 선택이므로 Luna의 4라운드 요청을 Astra의 4라운드 요청에 대한 반응으로 해석할 수 없다. 이 궤적만으로 억제 효과, 요청량의 확대, 혹은 tit-for-tat 정책을 확인할 수는 없다.

모델 호출 없이 \texttt{python3 arr2026/tools/}\allowbreak\texttt{extended\_analysis.py}로 분석을 재생성할 수 있다. 공개 점수, 칸별 횟수, 출처 링크는 \texttt{capability\_motive\_snapshot.csv}에, 원시 로그 경로와 줄 번호가 있는 에이전트·라운드 관측 449건은 \texttt{arena\_strategy\_observations.csv}에 보관했다. 8라운드 전체 후속 기록은 \texttt{astra\_luna\_followup.json}에 있다. 파일과 출처 해시는 \texttt{arr2026/results/} 아래에 있다. 단계별 분모의 차이를 유지했으며 반복 행동을 독립된 모델로 취급하지 않았다.
\FloatBarrier
